\section{COMPARISON OF FUNCTIONS IN A HARDY FIELD}

PROPOSITION 1. Two functions in the same Hardy field are comparable to any order (V, p.~232).

Indeed, if \(f\) belongs to a Hardy field \(\mathfrak{K}\) , then for every integer \(n > 0\) there exists an interval \([x_0, +\infty[\) on which \(f\) is \(n\) times differentiable, its \(n^{th}\) derivative being equal to a function in \(\mathfrak{K}\) on this interval. It is therefore enough to show that any two functions \(f, g\) of \(\mathfrak{K}\) are comparable. This is evident if one of the functions is identically zero on a neighbourhood of \(+\infty\) ; one may therefore restrict oneself to the case where they are both strictly positive on a neighbourhood of \(+\infty\) . But then, for every real number \(t\) , \(f - tg\) is equal to a function in \(\mathfrak{K}\) on a neighbourhood of \(+\infty\) , so has constant sign on a neighbourhood of \(+\infty\) , which proves the proposition (V, p.~217, prop. 9).

One deduces immediately from this proposition that, if a Hardy field \(\mathfrak{K}\) contains the real constants (as we shall always assume in what follows), and if \(f\) and \(g\) are any two functions in \(\mathfrak{K}\) then any two of the functions \(e^f\) , \(e^g\) , \(\log |f|\) , \(\log |g|\) , \(|f|^{\alpha}\) , \(|g|^{\alpha}\) ( \(\alpha\) an arbitrary real), \(\int_{a} f\) , \(\int_{a} g\) ( \(a\) being any real number in an interval \([x_0, +\infty[\) where \(f\) and \(g\) are regulated) are comparable (when they are defined); indeed, any two of these functions belong to a given Hardy field obtained by adjoining them successively to \(\mathfrak{K}\) .

Similarly, every function \(f(x)\) in a Hardy field \(\mathfrak{K}\) is comparable to \(x\) , since \(x\) and \(f(x)\) belong to the Hardy field obtained by adjoining \(x\) to \(\mathfrak{K}\) . One thus concludes (in particular) that \(f\) is comparable to any order to every power \(x^{\alpha}\) , as well as to \(\log x\) and to \(e^{x}\) .

One also sees that if \(f\) and \(g\) belong to the same Hardy field \(\mathfrak{K}\) , if \(g(x) > 0\) on an interval \([x_0, +\infty[\) , and if \(g(x)\) tends to 0 or to \(+\infty\) as \(x\) tends to \(+\infty\) , then the order of \(f\) relative to \(g\) (V, p.~219) is always defined.

Prop. 8 of V, p.~233, is therefore applicable to every function \(f\) in a Hardy field, and proves that:

\(1^{\circ}\) if \(f\) is of order \(+\infty\) relative to \(x\) then \(\int_{a}^{x} f(t) dt \sim (f(x))^{2} / f'(x)\) .

\(2^{\circ}\) if \(f\) is of order \(\mu > -1\) relative to \(x\) then \(\int_{a}^{x} f(t) dt \sim \frac{1}{\mu + 1} xf(x)\) .

\(3^{\circ}\) if \(f\) is of order \(\mu < -1\) relative to \(x\) then \(\int_{x}^{+\infty} f(t) dt \sim -\frac{1}{\mu + 1} xf(x)\) .

\(4^{\circ}\) if \(f\) is of order \(-\infty\) relative to \(x\) then \(\int_{x}^{+\infty} f(t) dt \sim -(f(x))^{2} / f'(x)\) .

Further, we have the following proposition:

PROPOSITION 2. Let \(f\) be a function belonging to a Hardy field \(\mathfrak{K}\) .

\(1^{\circ}\) If \(f\) is of infinite order relative to \(x\) , then, for every integer \(n > 0\) ,

\[
f ^ {(n)} (x) \sim \frac {(f ^ {\prime} (x)) ^ {n}}{(f (x)) ^ {n - 1}}.\tag{2}
\]

\(2^{\circ}\) If \(f\) is of finite order \(\mu\) relative to \(x\) then, for every \(n > 0\) ,

\[
\begin{array}{c} f ^ {(n)} (x) \sim \mu (\mu - 1) \ldots (\mu - n + 1) \frac {f (x)}{x ^ {n}} \\ \sim \frac {(\mu - 1) \ldots (\mu - n + 1)}{\mu^ {n - 1}} \frac {(f ^ {\prime} (x)) ^ {n}}{(f (x)) ^ {n - 1}} \end{array}\tag{3}
\]

except if \(\mu\) is an integer \(\geqslant 0\) and \(n > \mu\) .

\(1^{\circ}\) If \(f\) is of infinite order relative to \(x\) one has \(\log |f| \gg \log x\) , so, since \(\log |f|\) and \(\log x\) are comparable to any order, \(f'/f \gg 1/x\) . Put \(g = f'/f\) ; since \(g\) is equal to a function in \(\mathfrak{K}\) on a neighbourhood of \(+\infty\) one deduces from \(1/g \ll x\) , that \(g'/g^{2} \ll 1\) , and so \(g'/g \ll g = f'/f\) , or again \(fg' \ll gf'\) . From the relation \(f' = fg\) one deduces by differentiating that

\[
f ^ {\prime \prime} = f g ^ {\prime} + g f ^ {\prime} \sim g f ^ {\prime}
\]

or again \(f'' / f' \sim f' / f\) . The same argument, applied to \(f^{(n)}\) instead of \(f\) , shows, by induction on \(n\) , that \(f^{(n)} / f^{(n-1)} \sim f' / f\) ; whence relation (2).

\(2^{\circ}\) If \(f\) is of finite order \(\mu\) relative to \(x\) , and if \(\mu \neq 0\) , one has \(\log |f| \sim \mu \log x\) , whence, on differentiating, \(f'(x) \sim \mu \frac{f(x)}{x}\) ; one deduces that \(f'\) is of order \(\mu - 1\) relative to \(x\) , which allows one to apply the same argument by induction on \(n\) so long as \(\mu \neq n\) , whence formula (3) when \(\mu\) is not an integer \(\geqslant 0\) and \(< n\) .

When \(f\) is of integral order \(p \geqslant 0\) relative to \(x\) one can write \(f(x) = x^p f_1(x)\) , where \(f_1\) is of order 0 relative to \(x\) . By prop. 2 one has

\[
f ^ {(p)} \sim p! f _ {1}.
\]

To evaluate the derivatives of order \(n > p\) one may restrict oneself to the case \(p = 0\) . Then one has \(\log |f| \ll \log x\) , whence \(f'(x) / f(x) \ll 1 / x\) , in other words \(xf'(x) \ll f(x)\) ; if \(f\) is not equivalent to a constant \(k \neq 0\) one has, on differentiating this relation (V, p.~232, prop. 7), \(xf''(x) + f'(x) \ll f'(x)\) , which implies that \(xf''(x) \sim -f'(x)\) . Taking account of this formula, one sees by induction on \(n\) that \(f^{(n)}\) is of order \(\leqslant -n\) relative to \(x\) , and that

\[
f ^ {(n)} \sim (- 1) ^ {n + 1} (n - 1)! \frac {f ^ {\prime} (x)}{x ^ {n - 1}}.\tag{4}
\]

If \(f\) is equivalent to a constant \(k \neq 0\) one has \(f(x) = k + f_2(x)\) with \(f_2 \ll 1\) , and it reduces to studying the derivatives of \(f_2\) .

